Una gramática formal es un conjunto de reglas que definen cómo se construyen las oraciones válidas de un lenguaje. En lingüística computacional, las gramáticas libres de contexto (CFG) describen la estructura sintáctica de las oraciones mediante reglas de producción que descomponen cada oración en constituyentes (sujeto, verbo, complemento) y estos a su vez en palabras. A partir de estas reglas, un sistema puede generar automáticamente oraciones correctas dada una entrada estructurada.

En este proyecto, la gramática formal se aplica al problema de CAA: dada una secuencia de pictogramas seleccionados por el usuario, las reglas de producción determinan qué tipo de oración construir (afirmativa, interrogativa, negativa), conjugan los verbos, insertan los artículos y preposiciones necesarios, y producen una oración gramaticalmente correcta en español. La gramática se ha extendido con restricciones de selección semánticas que verifican la compatibilidad entre verbos y argumentos (por ejemplo, que ``comer'' requiere un objeto comestible) y con reglas de inferencia contextual que completan la oración cuando la entrada del usuario es incompleta.

\section{Planteamiento del enfoque}
\label{sec:f1-planteamiento}

La decisión de comenzar con una gramática formal responde a una motivación práctica: establecer un modelo base determinista donde la misma entrada produce siempre la misma salida, facilitando la evaluación. Una gramática ofrece control total sobre la generación, ya que cada transformación lingüística está codificada explícitamente, e independencia de dependencias externas.

La filosofía de diseño prioriza la corrección gramatical sobre la cobertura: en un contexto de CAA, donde el usuario depende del sistema para comunicarse, es preferible rechazar una entrada no soportada a generar una frase incorrecta que distorsione su intención comunicativa.

El sistema no emplea una gramática libre de contexto (CFG) pura en el sentido de Chomsky~\cite{chomsky1957syntactic}. En lugar de reglas de producción clásicas del tipo $S \rightarrow NP\ VP$, implementa una CFG extendida mediante tres mecanismos complementarios.

La gramática se define formalmente con Lark~\cite{lark2024}, un \textit{toolkit} de \textit{parsing} para Python que acepta gramáticas en notación EBNF. Las producciones principales describen la jerarquía de tipos oracionales y las estructuras sintagmáticas que el sistema reconoce.

\begin{lstlisting}[caption={Producciones principales de la gramática CFG (notación EBNF/Lark).},label={lst:cfg-producciones},basicstyle=\ttfamily\small,xleftmargin=1cm]
oracion: oracion_compuesta | oracion_simple

oracion_simple: oracion_interrogativa
              | oracion_negativa
              | oracion_afirmativa

oracion_afirmativa: oracion_necesidad | oracion_accion
                  | oracion_estado | oracion_descripcion
                  | oracion_locativa

oracion_necesidad: sujeto? verbo_necesidad complemento_necesidad
oracion_accion:    sujeto? verbo_accion complemento_verbal?
oracion_estado:    sujeto? verbo_estado predicado_estado
oracion_descripcion: sujeto verbo_copulativo predicado_descripcion
oracion_locativa:  sujeto? verbo_movimiento sintagma_preposicional

oracion_negativa:  NEGACION oracion_afirmativa

oracion_interrogativa: pregunta_con_interrogativo
                     | pregunta_si_no

oracion_compuesta: oracion_temporal | oracion_coordinada
oracion_temporal:  marcador_temporal oracion_simple
oracion_coordinada: oracion_simple conjuncion oracion_simple

sintagma_nominal:       determinante? sustantivo adjetivo?
sintagma_preposicional: preposicion sintagma_nominal
objeto_con_modificador: determinante? objeto adjetivo?
\end{lstlisting}

La gramática completa incluye además las definiciones de los terminales léxicos (pronombres, verbos por subcategoría, objetos, lugares, personas, adjetivos, interrogativos, marcadores temporales y modificadores), que en conjunto conforman las 646 entradas del léxico descrito en la sección~\ref{sec:f1-lexico}.

Lark permitió especificar estas reglas de forma declarativa y verificar su coherencia, pero el motor de generación no utiliza el \textit{parser} Earley ni el árbol de sintaxis que Lark produce. Las entradas de CAA son secuencias cortas (entre una y cinco palabras) extraídas de un vocabulario cerrado de 232 pictogramas, lo que hace que el conjunto de patrones válidos sea finito y enumerable. En este escenario, un \textit{parser} formal resulta sobredimensionado: la implementación traduce la gramática a 48 patrones sintácticos codificados como secuencias de categorías gramaticales y los reconoce por \textit{matching} directo. El resultado es más simple de depurar, determinista por construcción y suficiente para la complejidad del dominio.

Los tres mecanismos que extienden la CFG operan sobre estos patrones. El \textit{pattern-matching} sobre categorías gramaticales define cada patrón sintáctico como una secuencia ordenada de categorías (por ejemplo, verbo de necesidad seguido de verbo de acción y objeto), y se selecciona el primer patrón que coincide con la clasificación categorial de la entrada. Las restricciones de selección semánticas~\cite{resnik1996selectional} validan que los argumentos de cada verbo cumplan rasgos específicos (``comer'' requiere un objeto comestible, ``ir'' requiere un complemento de lugar). Las reglas de inferencia contextual completan la oración cuando la entrada del usuario es incompleta: por ejemplo, la entrada ``agua'' activa la inferencia del verbo ``quiero'', produciendo ``Quiero agua''.

Esta combinación permite manejar la ambigüedad inherente a la entrada abreviada de CAA sin recurrir a modelos estadísticos o neuronales.

\section{Decisiones de diseño específicas}
\label{sec:f1-postprocesamiento}

El postprocesador transforma la secuencia tokenizada en una oración natural. Tres de sus etapas incorporan decisiones de diseño específicas de este proyecto que van más allá de lo que una gramática formal estándar resolvería.

\subsection{Inferencia de elementos faltantes}

Esta etapa detecta elementos implícitos en la entrada del usuario y los añade explícitamente, siguiendo reglas que priorizan la inferencia conservadora:

\begin{itemize}
    \item Un solo objeto (``agua'') activa la inserción de ``quiero''.
    \item Un solo lugar (``baño'') activa la inserción de ``quiero'' e ``ir''.
    \item Un infinitivo sin verbo modal (``comer manzana'') recibe ``quiero'' como prefijo.
\end{itemize}

La inferencia opera exclusivamente en primera persona y no se activa en interrogativas, evitando asumir la intención del usuario cuando pregunta sobre terceros. Esta lógica es específica del dominio de CAA: en la comunicación aumentativa, el usuario selecciona los pictogramas mínimos necesarios y el sistema debe completar la intención comunicativa implícita.

\subsection{Inserción de determinantes}

La decisión de qué determinante acompaña a cada sustantivo depende de reglas específicas del dominio. Los sustantivos incontables (agua, pan, leche) no reciben determinante tras verbos de necesidad. Los objetos con artículo definido (televisión, teléfono, música) usan ``el/la'' en lugar de ``un/una''. En interrogativas se utiliza artículo definido (``¿Dónde está el perro?''). En coordinaciones, la coherencia entre los miembros determina el tratamiento: si un elemento es incontable, el otro también omite el determinante.

\subsection{Inserción de preposiciones}

La selección de la preposición correcta para complementos de lugar utiliza tres diccionarios construidos manualmente: preposiciones por verbo de movimiento (9 verbos: ``ir'' $\rightarrow$ ``a'', ``venir'' $\rightarrow$ ``de''), preposiciones por lugar con verbos de movimiento (35 lugares: ``baño'' $\rightarrow$ ``al'', ``escuela'' $\rightarrow$ ``a la'') y preposiciones de ubicación para verbos de estado (19 lugares: ``colegio'' $\rightarrow$ ``en el''). La lógica busca hacia atrás desde el lugar hasta encontrar el verbo que lo rige, y selecciona el diccionario apropiado según el tipo de verbo.

\section{Limitaciones}
\label{sec:f1-limitaciones}

Las gramáticas formales presentan limitaciones inherentes al enfoque simbólico. En primer lugar, operan sobre un vocabulario cerrado: toda palabra no incluida en el léxico se rechaza, lo que impide al sistema adaptarse a vocabulario nuevo sin intervención manual. En segundo lugar, las reglas de producción capturan la estructura sintáctica pero no el significado ni el contexto pragmático: una gramática no puede distinguir entre dos oraciones sintácticamente idénticas pero semánticamente diferentes, ni resolver ambigüedades que dependen del contexto comunicativo. En tercer lugar, la cobertura lingüística depende enteramente de las reglas codificadas: cada nueva estructura oracional, tiempo verbal o fenómeno morfológico requiere añadir reglas explícitamente, un esfuerzo de mantenimiento que crece con la complejidad del dominio. Por último, y quizá la limitación más relevante, las reglas son específicas de un idioma y un dominio concretos: las decisiones sobre inferencia de verbos modales, selección de preposiciones o inserción de determinantes están codificadas para el español y para el vocabulario de CAA de este proyecto. Extender el sistema a otro idioma o a un dominio comunicativo diferente requeriría reescribir la práctica totalidad de las reglas, ya que ninguna de ellas es transferible. Esta falta de portabilidad es inherente a las gramáticas formales aplicadas a la generación de lenguaje natural, y es precisamente lo que motiva y justifica la exploración de técnicas complementarias (embeddings semánticos, modelos de lenguaje) que puedan generalizar más allá del dominio para el que fueron diseñadas.
